\subsection{Reducción con Pila (ciclos en un bang path)}
\label{alg:3-12}

\graybold{Preguntas guía:}

\begin{itemize}
\item ¿Debo procesar una secuencia donde ciertos elementos pueden cancelarse al aparecer nuevamente?
\item ¿Eliminar un ciclo puede generar inmediatamente otro ciclo más grande?
\item ¿Necesito conservar el orden original mientras simplifico la secuencia al máximo?
\end{itemize}

\graybold{Utilidad:} Simplificar recorridos o secuencias eliminando ciclos innecesarios. Elimina dos rodeos locales: un elemento repetido dos veces seguidas (A A) y una ida y vuelta inmediata (A B A). Después de cada eliminación hay que volver a comprobar, porque pueden quedar juntos elementos que forman otro rodeo: así se deshacen en cascada rodeos anidados como a!b!c!b!a, pero no un ciclo general como a!b!c!a. Este patrón aparece en problemas de rutas, navegación, directorios, historial de movimientos y simplificación de secuencias.

\graybold{Complejidad:} \bigO{n}.

\graybold{Fundamento teórico:} La solución utiliza una pila que representa la ruta simplificada construida hasta el momento; cada máquina se procesa exactamente una vez y se inserta al final de la pila. Después de insertar una máquina pueden ocurrir dos situaciones. La primera es si las dos últimas máquinas tienen el mismo nombre donde la última visita es redundante y puede eliminarse. La segunda donde la última máquina coincide con la antepenúltima, por lo que tanto la máquina intermedia como la repetida pueden eliminarse.

\graybold{Ejercicio aplicado}

Dado un bang path de UUCP (UNIX-to-UNIX Copy Program), eliminar todos los saltos innecesarios producidos por ciclos entre máquinas, conservando exactamente el mismo destino final y respetando las mayúsculas y minúsculas de las máquinas que permanecen.

Las comparaciones entre nombres de máquinas son \textbf{case-insensitive}, aunque la salida debe conservar la capitalización original.

\graybold{I/O:} La entrada consiste en una única línea cuya longitud no supera los 256 caracteres. Cada componente está separado por el carácter '!' y el último elemento siempre corresponde al usuario destino, por lo que nunca participa en las simplificaciones. Debido al pequeño tamaño de la entrada, basta dividir el string utilizando \texttt{split("!")} y recorrer los componentes una única vez.

\graybold{Código solución}

\begin{lstlisting}[language=Python]
import sys
input = sys.stdin.buffer.readline
data = input().decode("utf-8").split("!") # Separa cada componente del bang path

if len(data) <= 2: # Si no existen maquinas intermedias no hay nada que simplificar
    sys.stdout.write("!".join(data))
    exit()

machines = data[:-1]     # todas las maquinas
user = data[-1]          # usuario final (no se modifica)
stack = []

for machine in machines:
    stack.append(machine)     # agrega la maquina actual a la ruta
    while True:     # mientras aparezcan nuevos ciclos despues de eliminar otros
        if (len(stack) >= 2 and stack[-1].lower() == stack[-2].lower()): #CASO AA
            stack.pop()
            continue

        if (len(stack) >= 3 and stack[-1].lower() == stack[-3].lower()): #CASO ABA
            stack.pop()      # elimina la ultima A
            stack.pop()      # elimina B
            continue
        break  # ya no existen reducciones posibles
sys.stdout.write("!".join(stack + [user])) # reconstruye el bang path conservando el usuario final
\end{lstlisting}
